\scanpage{224}
(\uline{N.B.} Ce dernier lemme est vrai pour $X/S$)

Soit en effet $F\in(x)\cap(y)\cap(x+y+z\tau)$, on a donc $F=xyG$, et $F$ devient nul en $y$ \unsure{faisant|\deleted{faisant}} $y=-(x+z\tau)$, d'où $x(x+z\tau)G(x,-x\unsure{+|-}z\tau,z)$ $=0$, or $x$ et $x+\tau z$ sont non diviseurs de $0$ dans $\underline{\mathcal O}_S[x,y,z]$ \inserted{$=\mathcal O_S[y,z][x]$} \deleted{\unsure{puisque}} d'où $G(x,-x-\tau z,z)=0$ i.e. $G$ \deleted{\unsure{est}} \inserted{devient} nul \deleted{par} la substitution $y=-x-z\tau$, donc $G$ est divisible par $y+x+z\tau$.

\uline{Corollaire du lemme 0} $\underline{\unsure{\mathcal J|\mathcal I|I}}_1\cap\underline{\mathcal J}_2\cap\underline{\mathcal J}_3=(0)$, donc
\[
\bigvee_{1\le i\le3}X_i=X.
\]

\uline{Lemme 2} \inserted{On a} \deleted{Si $i\ne j$, on a $\mathscr S/\unsure{\mathfrak p|p}_{i}+\mathfrak p_j\simeq\mathcal O_S[u]\unsure{[z]}$}
\[
\left\{\begin{aligned}
&\text{pr }i=1,2,3,\qquad S/\mathfrak p_i\simeq\underline{\mathcal O}_S[u,\unsure{v|o}]\qquad\text{donc }X_i\simeq P^1_S\\
&\text{Si }i\ne j,\qquad S/\mathfrak p_i+\mathfrak p_j\simeq\mathcal O_S[u]\quad\deleted{\unsure{\sim|\text{et}}}\ \text{donc }X_{ij}\unsure{\xrightarrow{\sim}|\simeq} S\\
&\deleted{\text{\unsure{Si |Enfin} }i}\qquad S/\mathfrak p_1+\mathfrak p_2+\mathfrak p_3\simeq\underline{\mathcal O}_S/\tau\underline{\mathcal O}_S[u]\quad\text{donc }X_{ijk}=(S,\mathcal O_S/\tau)\\
&&=\tau^{-1}(0)
\end{aligned}\right.
\]

\uline{Corollaire} Si $S$ est affine, on a
\[
H^1(X,\underline{\mathcal O}_X)\simeq H^1(X_1\amalg X_2\amalg X_3/X,\underline G_a)\simeq
\]
\inserted{justification : \deleted{\unsure{qccf}} $X_1\amalg X_2\amalg X_3\to X$ est un morphisme de descente \unsure{stricte} pour les \unsure{schémas}}
% archive sheet number: 223
